\documentclass[11pt]{article}

\usepackage[margin=1in]{geometry}
\usepackage{amsmath,amssymb}
\usepackage{array}
\usepackage{booktabs}
\usepackage{enumitem}
\usepackage[hidelinks]{hyperref}

\setlist[itemize]{leftmargin=1.4em}
\setlist[enumerate]{leftmargin=1.6em}

\title{SYSC 5108 Exam Study\\Assignment 4, Question 7}
\author{Jesse Levine}
\date{}

\begin{document}
\maketitle

\section*{Question}

We have the 1D data
\[
    x = [1,2,3,9,10].
\]
Assume a two-component Gaussian mixture model
\[
    p(x)=\pi_1 \mathcal{N}(x\mid \mu_1,\sigma_1^2)+\pi_2 \mathcal{N}(x\mid \mu_2,\sigma_2^2),
\]
with
\[
    \pi_1+\pi_2=1,
    \qquad
    \sigma_1^2=\sigma_2^2=1.
\]
The initial parameters are
\[
    \pi_1=0.6,\qquad \pi_2=0.4,\qquad \mu_1=2.0,\qquad \mu_2=9.5.
\]
Find the updated parameter values after one EM iteration.

\section*{Course and Textbook Cross-References}

\begin{itemize}
    \item \textbf{Chapter 7 slides, Gaussian Mixture Model:} the data density is a
    weighted sum of Gaussian components with mixing coefficients summing to 1.
    \item \textbf{Chapter 7 slides, EM Algorithm:} first compute the hidden-variable
    responsibilities in the E-step, then update mixture weights and means in the
    M-step.
    \item \textbf{Goodfellow, Bengio, and Courville, \emph{Deep Learning}, mixture models:}
    Gaussian mixtures introduce latent component assignments.
    \item \textbf{Goodfellow, Bengio, and Courville, \emph{Deep Learning}, EM discussion:}
    EM alternates between estimating latent-variable posteriors and maximizing the
    expected complete-data log-likelihood.
\end{itemize}

\section*{Step 1: Write the Gaussian Density}

Because the variance is fixed at \(1\), each component density is
\[
    \mathcal{N}(x\mid \mu,1)
    =
    \frac{1}{\sqrt{2\pi}}
    \exp\!\left(-\frac{(x-\mu)^2}{2}\right).
\]

In the E-step, the common factor \(1/\sqrt{2\pi}\) cancels in the responsibility
ratio, so we only need the exponential terms.

\section*{Step 2: E-Step Responsibilities}

Let \(\gamma_{i1}\) be the responsibility of component \(1\) for point \(x_i\). Then
\[
    \gamma_{i1}
    =
    \frac{\pi_1 \mathcal{N}(x_i\mid \mu_1,1)}
    {\pi_1 \mathcal{N}(x_i\mid \mu_1,1)+\pi_2 \mathcal{N}(x_i\mid \mu_2,1)},
\]
and
\[
    \gamma_{i2}=1-\gamma_{i1}.
\]

Using
\[
    (\pi_1,\pi_2,\mu_1,\mu_2)=(0.6,0.4,2.0,9.5),
\]
we evaluate the responsibilities for the five data points.

\subsection*{For \(x_1=1\)}

\[
    \gamma_{11}
    =
    \frac{0.6\,e^{-(1-2)^2/2}}
    {0.6\,e^{-(1-2)^2/2}+0.4\,e^{-(1-9.5)^2/2}}
    =
    \frac{0.6e^{-0.5}}{0.6e^{-0.5}+0.4e^{-36.125}}
    \approx 1.
\]
So
\[
    \gamma_{11}\approx 1,\qquad \gamma_{12}\approx 0.
\]

\subsection*{For \(x_2=2\)}

\[
    \gamma_{21}
    =
    \frac{0.6\,e^{-(2-2)^2/2}}
    {0.6\,e^{-(2-2)^2/2}+0.4\,e^{-(2-9.5)^2/2}}
    =
    \frac{0.6}{0.6+0.4e^{-28.125}}
    \approx 1.
\]
So
\[
    \gamma_{21}\approx 1,\qquad \gamma_{22}\approx 0.
\]

\subsection*{For \(x_3=3\)}

\[
    \gamma_{31}
    =
    \frac{0.6\,e^{-(3-2)^2/2}}
    {0.6\,e^{-(3-2)^2/2}+0.4\,e^{-(3-9.5)^2/2}}
    =
    \frac{0.6e^{-0.5}}{0.6e^{-0.5}+0.4e^{-21.125}}
    \approx 1.
\]
So
\[
    \gamma_{31}\approx 1,\qquad \gamma_{32}\approx 0.
\]

\subsection*{For \(x_4=9\)}

\[
    \gamma_{41}
    =
    \frac{0.6\,e^{-(9-2)^2/2}}
    {0.6\,e^{-(9-2)^2/2}+0.4\,e^{-(9-9.5)^2/2}}
    =
    \frac{0.6e^{-24.5}}{0.6e^{-24.5}+0.4e^{-0.125}}
    \approx 0.
\]
So
\[
    \gamma_{41}\approx 0,\qquad \gamma_{42}\approx 1.
\]

\subsection*{For \(x_5=10\)}

\[
    \gamma_{51}
    =
    \frac{0.6\,e^{-(10-2)^2/2}}
    {0.6\,e^{-(10-2)^2/2}+0.4\,e^{-(10-9.5)^2/2}}
    =
    \frac{0.6e^{-32}}{0.6e^{-32}+0.4e^{-0.125}}
    \approx 0.
\]
So
\[
    \gamma_{51}\approx 0,\qquad \gamma_{52}\approx 1.
\]

\subsection*{Responsibility Table}

To four decimal places:

\[
\begin{array}{c|cc}
    x_i & \gamma_{i1} & \gamma_{i2}\\
    \hline
    1  & 1.0000 & 0.0000\\
    2  & 1.0000 & 0.0000\\
    3  & 1.0000 & 0.0000\\
    9  & 0.0000 & 1.0000\\
    10 & 0.0000 & 1.0000
\end{array}
\]

Hence the effective component counts are
\[
    N_1=\sum_{i=1}^5 \gamma_{i1}\approx 3,
    \qquad
    N_2=\sum_{i=1}^5 \gamma_{i2}\approx 2.
\]

\section*{Step 3: M-Step Updates}

With fixed variance, the updated mixing coefficients and means are
\[
    \pi_1^{\text{new}}=\frac{N_1}{5},
    \qquad
    \pi_2^{\text{new}}=\frac{N_2}{5},
\]
\[
    \mu_1^{\text{new}}
    =
    \frac{\sum_{i=1}^5 \gamma_{i1}x_i}{N_1},
    \qquad
    \mu_2^{\text{new}}
    =
    \frac{\sum_{i=1}^5 \gamma_{i2}x_i}{N_2}.
\]

\subsection*{Update the Mixing Coefficients}

\[
    \pi_1^{\text{new}}=\frac{3}{5}=0.6,
    \qquad
    \pi_2^{\text{new}}=\frac{2}{5}=0.4.
\]

\subsection*{Update the Means}

For component \(1\),
\[
    \mu_1^{\text{new}}
    =
    \frac{1+2+3}{3}
    =
    2.0.
\]

For component \(2\),
\[
    \mu_2^{\text{new}}
    =
    \frac{9+10}{2}
    =
    9.5.
\]

\section*{Final Answer}

After one EM iteration,
\[
    \boxed{\pi_1^{\text{new}}=0.6000,\qquad \pi_2^{\text{new}}=0.4000,}
\]
\[
    \boxed{\mu_1^{\text{new}}=2.0000,\qquad \mu_2^{\text{new}}=9.5000.}
\]

\section*{Why Nothing Changed}

The data are already almost perfectly separated into two groups:
\[
    \{1,2,3\}
    \qquad \text{and} \qquad
    \{9,10\}.
\]
The initial means \(2.0\) and \(9.5\) already sit at the centers of those groups,
and the initial weights \(0.6\) and \(0.4\) already match the proportions
\[
    \frac{3}{5}
    \qquad \text{and} \qquad
    \frac{2}{5}.
\]
So the E-step assigns almost all responsibility to the obvious component for each
point, and the M-step returns essentially the same parameters.

\section*{Exam Shortcut}

For a clearly separated two-component GMM question:
\begin{enumerate}[label=\arabic*.]
    \item Compute the responsibilities \(\gamma_{ik}\).
    \item Sum them to get \(N_k\).
    \item Update \(\pi_k=N_k/M\).
    \item Update each mean by the responsibility-weighted average.
\end{enumerate}

If the clusters are far apart, the responsibilities will be nearly \(0\) or \(1\),
which makes the arithmetic very fast.

\section*{Sources Used}

\begin{itemize}
    \item Assignment 4 PDF, Question 7.
    \item Local submitted Assignment 4 PDF checked for consistency of the final values.
    \item Chapter 7 slides on Gaussian mixture models and the EM algorithm.
    \item Goodfellow, Bengio, and Courville, \emph{Deep Learning}, sections on mixture models and EM.
\end{itemize}

\end{document}
